\section{From Euclidean rotations to Lorentz covariance}
\label{sec:lorentz-covariance}
This section proves covariance of the boundary distributions from the
Euclidean rotation identities, retaining the same field labels. The
finite-dimensional Lorentz action on those labels must first be available.
We state its required compatibility explicitly rather than identify it with
an arbitrary representation having the same dimension.

On each finite level let $R$ be the Euclidean rotation representation, and
write $L_a$ for its infinitesimal generator in the plane consisting of time
and spatial coordinate $a$. Our convention for the coordinate generator is
\[
 \mathcal R_a=\sum_i\left(x_i^a\partial_{y_i}
                                  -y_i\partial_{x_i^a}\right).
\]
Assume $R$ acts compatibly under the inclusions of field levels, and the
smooth Euclidean correlation $S(v;X)$ obeys simultaneous rotation covariance.
Differentiating it gives
\begin{equation}\label{eq:euclid-infinitesimal}
 \mathcal R_aS(v;X)=-\sum_i S(v_1,\ldots,L_av_i,\ldots,v_n;X).
\end{equation}
Indeed, covariance says $S(v;R_tX)=S(R_{-t}v;X)$; the ordinary chain rule on
positions and the finite multilinear expansion on labels give the displayed
formula. This is applied only at distinct Euclidean positions.

Suppose also that the proper orthochronous Lorentz group has a representation
$\rho$ on the field union preserving every finite level, continuous there,
compatible with the inclusions, and agreeing with $R$ on spatial rotations.
For the coordinate boost $B_a(s)$, with coordinate infinitesimal generator
$x^a\partial_t+t\partial_{x^a}$, require
\begin{equation}\label{eq:label-boost-input}
 \left.\frac{d}{ds}\rho(B_a(s))v\right|_{s=0}=-iL_av.
\end{equation}
Proposition~\ref{prop:lorentz-label-descent} supplies this action from the
calibrated finite-level actions. Their construction remains an explicit
earlier input. Set $K_a=-iL_a$.

\begin{proposition}[The tested boost identity]\label{prop:boost-ward}
For strictly ordered $y$ and any joint Schwartz test $f$,
\begin{equation}\label{eq:boost-interior}
 \sum_i\psi_y(v_1,\ldots,K_av_i,\ldots,v_n;f)
  =\psi_y\left(v;\sum_i
       [x_i^a\partial_{t_i}+(t_i+iy_i)\partial_{x_i^a}]f\right).
\end{equation}
Consequently the boundary distributions satisfy
\begin{equation}\label{eq:boost-boundary}
 \sum_i W_n(v_1,\ldots,K_av_i,\ldots,v_n)(f)
  =W_n(v)(\mathcal B_af),\qquad
 \mathcal B_a=\sum_i(x_i^a\partial_{t_i}+t_i\partial_{x_i^a}).
\end{equation}
\end{proposition}
\begin{proof}
First suppose $f$ has compact support $C$. For a complex parameter $c$ put
$z_i(c)=iy_i+ct_i$ and
\[
 F_c(v;X)=F_n(z(c),x;v),\qquad
 T_c f=\sum_i[x_i^a\partial_{t_i}+c z_i(c)\partial_{x_i^a}]f.
\]
Let $D$ consist of parameters for which $z(c)$ has strictly ordered
imaginary parts for every $X\in C$. This is an open convex set containing
the real axis: its constraints are affine strict inequalities in
$\operatorname{Im}c$, and compactness of $C$ makes their margins uniform
near each admissible parameter. Consider the integral
\[
 H(c)=\int_C\left[F_c(v;X)T_cf(X)
       -c\sum_iF_c(v_1,\ldots,K_av_i,\ldots,v_n;X)f(X)\right]dX.
\]
It is holomorphic on $D$. In fact the integrand is jointly continuous on
$D\times C$ and holomorphic in $c$ for every $X$; on each compact subdisc,
Cauchy's formula gives uniformly bounded derivatives, so differentiation
under this finite-volume integral is justified.

At $c=i\sigma$, put $Y_i=y_i+\sigma t_i$. All these heights are strictly
ordered on a neighborhood of $C$, and $F_{i\sigma}(v;X)=S(v;(x,Y))$.
The real vector field
\[
 A_\sigma=\sum_i\left[x_i^a\partial_{t_i}
                      -\sigma(y_i+\sigma t_i)\partial_{x_i^a}\right]
\]
is divergence free. Under the map $(t,x)\mapsto(Y,x)$ it becomes
$\sigma\mathcal R_a$. Thus \eqref{eq:euclid-infinitesimal} gives
\[
 A_\sigma F_{i\sigma}(v;X)
       =-\sigma\sum_iS(v_1,\ldots,L_av_i,\ldots,v_n;(x,Y))
       =-i\sigma\sum_i F_{i\sigma}(v_1,\ldots,K_av_i,\ldots,v_n;X).
\]
Also $T_{i\sigma}f=A_\sigma f$. The integrand of $H(i\sigma)$ is
therefore $A_\sigma(F_{i\sigma}(v;X)f(X))$, whose integral is zero.
The product is smooth near $C$ and zero off $C$, so its extension by zero
is a smooth compactly supported function; no regularity of the continued
function at positions outside that neighborhood is being assumed.
The identity theorem on the connected set $D$ now gives $H(1)=0$,
which is \eqref{eq:boost-interior} for compactly supported tests.

To pass to Schwartz tests, choose a smooth cutoff $\chi(X/R)$ equal to
one near zero. Then $\chi(X/R)f\to f$ in every Schwartz seminorm: on the
outside region use one extra polynomial decay weight, and for each
positive-order derivative of the cutoff use its factor $R^{-1}$ or a
higher inverse power and the same tail estimate. At fixed strictly ordered
$y$, every $\psi_y$ is continuous on Schwartz space, and all coordinate
multiplications and derivatives in \eqref{eq:boost-interior} are continuous
there. The compact-test identity consequently extends to $f$.

Finally take the full ordered-height limit $y\to0$. The test on the
right is $\mathcal B_af+\sum_i iy_i\partial_{x_i^a}f$. Each fixed-test
pairing $\psi_y(v;\partial_{x_i^a}f)$ has a finite boundary limit, so
multiplication by $y_i$ makes the correction tend to zero. All other terms
have their specified joint boundary limits, proving
\eqref{eq:boost-boundary}.
\end{proof}

\begin{lemma}[Affine pullback of tests]\label{lem:schwartz-motion}
The action $P_gf(X)=f(g^{-1}X)$ of the Poincar\'e group, acting on all
positions at once, is continuous in the Schwartz topology. For coordinate
boosts,
\[
 \left.\frac{d}{ds}P_{B_a(s)}f\right|_{s=0}=-\mathcal B_af
\]
in that topology.
\end{lemma}
\begin{proof}
On a compact neighborhood of a fixed invertible affine map, its linear
part and inverse have uniformly bounded norms and its translation is
bounded. Hence $1+|X|$ and $1+|g^{-1}X|$ are comparable uniformly there.
The chain rule expresses every derivative of $f\circ g^{-1}$ as a finite
sum of derivatives of $f$ at $g^{-1}X$ times products of matrix entries.
Each weighted supremum seminorm is therefore bounded by finitely many
Schwartz seminorms of $f$, uniformly on that neighborhood.

For continuity, split the supremum into a fixed ball and its complement.
On the ball, uniform continuity of each derivative and continuous variation
of the matrix coefficients give convergence. On the complement, the
comparison of weights and one extra decay seminorm make the tail uniformly
small. This proves continuity in every seminorm.
For the boost derivative, the pointwise chain rule at zero gives
$-\mathcal B_af$. To justify the same derivative in Schwartz space, use
Taylor's formula along the smooth curve of invertible matrices
$B_a(-s)$. Its second derivative acts on $f$ as a finite sum of first and
second derivatives composed with $B_a(-s)$, multiplied by coefficients
polynomial in $X$ of degree at most two. The preceding weight comparison
bounds every seminorm of this second derivative uniformly for $s$ in a
small compact interval. The integral Taylor remainder is therefore
$O(s^2)$ in each seminorm, proving the stated derivative.
\end{proof}

\begin{theorem}\label{thm:lorentz-covariance}
Under the Euclidean covariance and compatible label-action hypotheses above,
the boundary distributions satisfy
\[
 W_n(\rho(\Lambda)v_1,\ldots,\rho(\Lambda)v_n)(P_\Lambda f)
          =W_n(v)(f)
\]
for every proper orthochronous Lorentz transformation $\Lambda$.
They are also invariant under common spacetime translations.
\end{theorem}
\begin{proof}
For a fixed coordinate boost set
\[
 A(s)=W_n(\rho(B_a(s))v_1,\ldots,\rho(B_a(s))v_n)
                                           (P_{B_a(s)}f).
\]
Put the finitely many labels in one common finite level. Expanding them in
a finite basis expresses $A(s)$ as a finite sum of smooth scalar
coefficients times fixed-label distributions applied to the moving test.
The label coefficients are smooth along the boost: its group law and
the derivative at zero give the linear differential equation
$\rho(B_a(s))' = \rho(B_a(s))K_a$ on each finite level.
Thus Lemma~\ref{lem:schwartz-motion} and
\eqref{eq:label-boost-input} justify the ordinary product rule at zero.
Its label terms are the left side of \eqref{eq:boost-boundary}, and its
test term is the negative of the right side. Hence $A'(0)=0$.
The representation laws for labels and pullbacks give the same calculation
at any $s$, using the labels and test already moved by $B_a(s)$.
Therefore $A'(s)=0$ everywhere and $A(s)=A(0)$.

Spatial rotations and common spatial translations are inherited by the
strict-time continuation from their Euclidean identities. To check this
without an assumption of holomorphy in spatial coordinates, fix all real
positions, compare the two scalar continued expressions, and replace each
time by $iy_i+ct_i$. Both sides are holomorphic on an open strip about the
real axis and agree for $c=i\sigma$ near zero by Euclidean covariance.
The identity theorem gives equality at $c=1$.
For a common real-time translation $b$, compare times
$iy_i+c(t_i+b)$ on one side with $iy_i+ct_i$ on the other. Their gaps
agree, and at $c=i\sigma$ they differ by the common Euclidean time
translation $\sigma b$. The same proof applies. Changing variables in the
joint test integrals, with unit absolute Jacobian, and taking boundary
limits gives spatial-rotation covariance and translation invariance of $W$.

Finally these spatial rotations and coordinate boosts generate the proper
orthochronous Lorentz group. Here is a direct verification. Given
$\Lambda$, its image $q=\Lambda e_0$ is future timelike with Minkowski
norm squared one. A boost in coordinate $a$ with rapidity
$-\operatorname{artanh}(q^a/q^0)$ makes that spatial component zero;
$|q^a|<q^0$ ensures the parameter exists. It leaves the other spatial
components unchanged. Three such boosts produce a vector with zero spatial
part and positive time component, necessarily $e_0$. Their product $B$
therefore makes $B\Lambda$ fix $e_0$. Preservation of the Minkowski form
then shows that $B\Lambda$ preserves its spatial orthogonal complement
and is orthogonal on it; determinant one makes its spatial part a rotation.
Thus $\Lambda$ is a product of coordinate boosts and a spatial rotation.
Covariance is closed under products and inverses by the representation
laws, proving the assertion for every $\Lambda$.
\end{proof}

Together with Corollary~\ref{cor:cone-support}, this gives the full forward
suffix-cone support of the boundary distributions, conditional only on the
specified compatible Lorentz action and the previously stated Euclidean
construction inputs. It does not yet prove spacelike commutativity.
